\section{Algorithm Details}
\label{sec:appendix_algorithms}

This appendix provides the full mathematical specifications for all algorithm variants summarised in the main paper.

\subsection{Synchronous Q-Learning}

In \emph{synchronous} Q-learning, the agent updates the Q-values for \emph{all} actions from the same state. Let $A$ be the set of possible actions. If the agent took action $a_t$ but also computes hypothetical rewards $r_t(a)$ for each $a \in A$, the synchronous update is:
\begin{equation}
Q\bigl(s_t,a\bigr)
\;\leftarrow\;
\bigl(1-\alpha\bigr)\,Q\bigl(s_t,a\bigr)
\;+\;
\alpha \Bigl[
  r_t(a)
  \;+\;
  \gamma \,\max_{a'} Q\bigl(s_{t+1},a'\bigr)
\Bigr]
\quad
\forall a \in A.
\label{eq:qlearning_sync}
\end{equation}
Thus every $Q(s_t,a)$ is updated in a single step, including actions the agent did not actually choose.

\subsection{Exploration Strategies}

\subsubsection{Boltzmann Exploration}

\emph{Boltzmann} (or \emph{softmax}) exploration draws an action $a$ from a distribution favouring higher estimated $Q$-values:
\begin{equation}
P\bigl(a \mid s\bigr)
\;=\;
\frac{\exp\!\bigl(\beta\,Q(s,a)\bigr)}{\sum_{b}\,\exp\!\bigl(\beta\,Q(s,b)\bigr)},
\label{eq:boltzmann}
\end{equation}
where $\beta = 1/\tau$ is the inverse temperature parameter. Larger $\beta$ (lower temperature $\tau$) makes the distribution more peaked, favouring higher-valued actions; smaller $\beta$ (higher $\tau$) flattens the distribution, promoting more exploration.

\subsubsection{Epsilon-Greedy Exploration}

\emph{$\varepsilon$-greedy} exploration selects the action that maximises $Q(s,a)$ with probability $1-\varepsilon$ and picks an action uniformly at random with probability $\varepsilon$:
\begin{equation}
P\bigl(a \mid s\bigr)
\;=\;
\begin{cases}
1-\varepsilon, & \text{if }a = \displaystyle\arg\max_{b}\,Q(s,b),\\
\displaystyle\frac{\varepsilon}{|A|}, & \text{otherwise}.
\end{cases}
\label{eq:egreedy}
\end{equation}
As $\varepsilon$ diminishes, the policy exploits more aggressively based on current $Q$-estimates.

\subsection{Multi-Armed Bandits and UCB}

A \emph{multi-armed bandit} scenario omits state transitions, focusing instead on learning which of several actions (arms) maximises expected reward. Let $\hat{\mu}_a$ be the estimated mean reward of arm $a$, and $u_a$ be an uncertainty term. A typical approach is Upper Confidence Bound (UCB), which selects
\begin{equation}
a_t
\;=\;
\arg\max_{a}\Bigl[\hat{\mu}_a + u_a\Bigr],
\end{equation}
then updates $\hat{\mu}_a$ and $u_a$ using the reward observed after pulling arm $a$. In our implementation, we initialise $\mathbf{A}_a = \lambda \mathbf{I}$ and $\mathbf{b}_a = \mathbf{0}$ for each arm $a$. The regularisation parameter $\lambda$ ensures that $\mathbf{A}_a$ is invertible from the start and provides numerical stability.

\subsection{Contextual Thompson Sampling}

As an alternative to the optimism-based exploration of LinUCB, \emph{Contextual Thompson Sampling} (CTS) uses posterior sampling to balance exploration and exploitation. Each action $a$ maintains a Bayesian linear regression model with posterior $\boldsymbol{\theta}_a \mid \text{data} \sim \mathcal{N}(\hat{\boldsymbol{\theta}}_a,\, \sigma^2 \mathbf{A}_a^{-1})$, where $\hat{\boldsymbol{\theta}}_a = \mathbf{A}_a^{-1}\mathbf{b}_a$ and $\sigma^2$ is a noise variance parameter. At each step, the algorithm draws a sample $\tilde{\boldsymbol{\theta}}_a$ from the posterior for every action and selects greedily with respect to the samples:
\begin{equation}
a_t \;=\; \arg\max_{a}\; \tilde{\boldsymbol{\theta}}_a^\top \mathbf{x}, \quad \tilde{\boldsymbol{\theta}}_a \sim \mathcal{N}\!\bigl(\hat{\boldsymbol{\theta}}_a,\, \sigma^2 \mathbf{A}_a^{-1}\bigr).
\end{equation}
The parameter $\sigma^2$ plays an analogous role to $c$ in LinUCB: larger $\sigma^2$ produces wider posterior draws, encouraging more exploration; smaller values exploit the current estimates more aggressively.

\subsection{Implementation Notes}

Each experiment runs for a set number of \emph{episodes}, during which agents collect rewards and update their parameters. Both $\varepsilon$-greedy and Boltzmann exploration begin with high randomness ($\varepsilon = 1.0$ or temperature $\tau = 1.0$) and decay linearly or exponentially across the first 90\% of episodes, settling to a floor of 0.01. In the final 10\% of episodes, $\varepsilon$-greedy agents switch to pure exploitation ($\varepsilon = 0$) and Boltzmann agents use a near-greedy temperature ($\tau = 0.01$). Q-tables may be initialised to all zeros or to small random values, influencing early exploration. In any given state, if multiple actions share the same $Q$-value, ties are broken uniformly at random. In the bandit setting, particularly with \emph{LinUCB}, two key parameters shape the agent's behaviour: the regularisation parameter $\lambda$, which scales the identity matrix added to the design matrix $\mathbf{X}^\top \mathbf{X}$ to stabilise estimates, and the exploration coefficient $c$, which controls the width of the confidence interval around each action's estimated reward. Larger $c$ promotes more exploration, whereas smaller $c$ favours exploiting the current reward estimates.

\subsection{PID Pacing}

This controller is used in Experiment~4b; Experiment~4a uses the multiplicative pacing described in the main paper. As an alternative to the multiplicative dual update of Balseiro and Gur (2019), we implement a proportional--integral (PI) controller that operates directly on the cumulative spending error:
\begin{align}
e_t^i &\;=\; \tfrac{t}{T}\,B^i \;-\; {\textstyle\sum_{\tau \le t}} c_\tau^i, \label{eq:pid_error}\\
\lambda_{t+1}^i &\;=\; \operatorname{clip}\!\bigl(\lambda_t^i + K_P\,e_t^i + K_I\!\textstyle\sum_{\tau} e_\tau^i,\; 0.01,\; 1.5\bigr), \label{eq:pid_update}
\end{align}
where $K_P = 0.30\times\text{aggressiveness}$ and $K_I = 0.05\times\text{aggressiveness}$. A positive error (underspending) raises $\lambda$, encouraging higher bids; overspending drives $\lambda$ down. The derivative term is omitted ($K_D = 0$) following the literature consensus that it is counterproductive in stochastic auction environments.
